


%\hed{Coloring and recoloring}
\section{Coloring and recoloring}


The rank condition {\bf not} holding means that for every coloring $\phi$, there is a
polynomial $p_\phi$, of degree at most $d$, with $p_\phi(0)=0$,
such that $\Omega p_\phi(\phi\Dl)=0$.



Introduce variables $\vep_i$, $\vep_i^2=1$. Then generically
$$\phi=\hbox{diag}(\vep_1,\vep_2,\ldots,\vep_n)=
  \begin{pmatrix} 
	\vep_1&0&\cdots&0\cr 0&\vep_2&\cdots&0  \\ %\cr
          0&0&\ddots&0 \\		%\cr
          0&0&\cdots&\vep_n
   \end{pmatrix}	%\cr}
$$


We want the matrices $\{\Om\Delta,\,\Om\Delta^2,\,\ldots, \Om\Delta^d\}$ to be linearly
independent. Since these are all KROSPS, we form a matrix $H$ whose $i^{\rm th}$ row
is $\pi(\phi\Dl)^i$. Then the rank of $H$ is $d$ if and only if $\det HH^*\ne0$.
\bigskip


\begin{remark} Note that $\Om$ is always linearly independent of the $\Om\Dl^i$, just apply
to the all 1's vector. So if $d=n-1$, you can include $\pi$ as the first row of $H$ and get
an $n\times n$ matrix. Then you only need to check $\det H$.
\end{remark}



\subsection*{Procedure}

Initialize $\phi_0=\hbox{diag}(1,1,\ldots,1)$. 

\begin{enumerate}
\item[1.] Form $HH^*$.
\item[2.] Calculate $D(\phi)=\det HH^*$.
\item[3.] If $D(\phi_0)\ne0$, {\bf DONE} else continue.
\item[4.] Compute $\nabla_\vep D(\phi_0)$.
\item[5.] If any component of $\nabla_\vep D(\phi_0)$ is nonzero, recolor the corresponding vertex, {\bf DONE} else continue.
\item[6.] If all components of the gradient are zero, randomly pick a new $\phi_0$. Repeat from Step 3.
\end{enumerate}




% CC existence criterion added in the development manuscript.
